\documentclass{article}
\usepackage{xfp}
\usepackage{tikz}
\usetikzlibrary{math}

\input{../../tikzlibrarymc.code.tex}

\begin{document}

% -----------------------------------------------
% \tikzmath{
%   print{\ \\===== Linear Equations =====};
%   \a{1,1} = 2; \a{1,2} = 1; \a{1,3} = 3; \a{1,4} = 4; 
%   \a{2,1} = 1; \a{2,2} = -2; \a{2,3} = 2; \a{2,4} = 3; 
%   \a{3,1} = 3; \a{3,2} = 4; \a{3,3} = -1; \a{3,4} = 2; 
%   \a{4,1} = 4; \a{4,2} = 3; \a{4,3} = 2; \a{4,4} = 1;
%   \b{1,1} = 16; \b{1,2} = 20;
%   \b{2,1} = 19; \b{2,2} = 11;
%   \b{3,1} = -1; \b{3,2} = -3;
%   \b{4,1} = 4; \b{4,2} = 5;
% }

% \tikzset{
%   mat/solve={\a,\b,4,2,\c},
%   mat/show={\a,4,4},
%   mat/show={\b,4,2},
%   mat/show={\c,4,2},% 1,2,-1,3
% }

% \tikzmath{
%   print{\ \\status:\status};
% }

% -----------------------------------------------
\tikzmath{
  print{\ \\===== Inverse =====};
  \a{1,1} = 2; \a{1,2} = 3; \a{1,3} = 4; \a{1,4} = 5; 
  \a{2,1} = 9; \a{2,2} = 3; \a{2,3} = 1; \a{2,4} = 2; 
  \a{3,1} = 4; \a{3,2} = 7; \a{3,3} = 6; \a{3,4} = 3; 
  \a{4,1} = 5; \a{4,2} = 9; \a{4,3} = 8; \a{4,4} = 7;
}

\tikzset{
  mat/inverse={\a,4,\b},
  mat/product={\a,\b,4,4,4,\c},
  mat/show={\a,4,4},
  mat/show={\b,4,4},
  mat/show={\c,4,4},
}

\tikzmath{
  print{\ \\status:\status};
}

\end{document}